\section{Introduction}\label{sec:intro}

Most DC--DC converters in industrial equipment are still closed by a PID-type
loop. The reason is practical: a PID block is available in every digital
controller, technicians know how to read its three gains, and it needs a single
voltage sensor. Model-based laws such as the linear--quadratic regulator are
better understood mathematically, but they ask for more sensors and their
tuning knobs, the weights of a quadratic cost, have no obvious meaning on the
shop floor. ISA Transactions addresses the engineering of measurement and
automation for practitioners and applied researchers, and this gap between
what theory certifies and what industry implements is the subject of this
paper. We study one converter in detail, a non-ideal asynchronous Buck stage fed
from \SI{48}{V} and commanded from \SI{1}{V} to \SI{46}{V}, and ask a precise
question: can a controller be tuned in the coordinates of optimal control and
still be delivered to the plant as an ordinary PID with a derivative filter,
without losing the optimality certificate on the way?

The answer turns out to be yes, and exactly so, under conditions that we state
and test. The rest of this section reviews how PID-type and linear--quadratic
controllers reached their present form and why the non-ideal Buck is a
demanding test case. It then states three research gaps, each paired with the
design choice it motivates (the PID equivalent, the CMO--LQI formulation, the
metaheuristics and refineBO), followed by the working assumptions, the
difficulties and the contributions.

\subsection{Literature review}\label{sec:lit}

\subsubsection{From the three-term controller to the 2-DOF PIDF}
The three-term law was first analysed by \citet{Minorsky1922} for the automatic
steering of ships, and the rules published by \citet{Ziegler1942} made it an
industrial standard. On a Buck converter the reason for each term is easy to
see. With proportional action alone the output settles below its reference
whenever the load draws current, because only a standing error can pull the
duty ratio away from its feedforward value. Integral action removes that
offset, which gives the PI controller. The output filter, however, is lightly
damped (in our model $\zeta\approx0.14$ at \SI{24}{V}), and a PI has no phase
lead to offer it; the derivative term supplies that lead and turns the PI into
a PID \citep{Astrom2006}. An ideal derivative cannot be built, and on a
converter it would also amplify the switching ripple riding on the measured
voltage. Industrial controllers and the Simulink PID block therefore filter it
with a first-order pole $N$, and the PID becomes a PIDF
\citep{Ang2005,Astrom2006}.

The PIDF still uses one error signal for two different jobs. The gains that
make it recover quickly from a load step act in the same way on every change
of reference, so a loop tuned for load rejection tends to overshoot on a large
setpoint step, such as the \SI{1}{}$\to$\SI{46}{V} step of our mission.
Setpoint weighting decouples the two. The proportional and derivative terms act
on $b\,r-y$ and $c\,r-y$ instead of $r-y$, while the integral keeps acting on the
true error, so the steady state does not change \citep{Astrom2006,Astrom2001};
\citet{Bianchi2008} extended the idea to multivariable PID loops. The general
principle goes back to \citet{Horowitz1963}: one path of the controller is
designed for robustness and disturbance rejection, and a second path, driven by
the reference alone, shapes the command response. Setpoint weighting is the
simplest instance of this two-degree-of-freedom (2-DOF) structure. Adding $N$,
$b$ and $c$ to the three gains gives a law with six parameters,
$(K_p,K_i,K_d,N,b,c)$, which is as far as a standard PID block goes; we take it
as the end point of the upper branch of Fig.~\ref{fig:evolution}. Surveys of
PID practice document both the dominance of this family in industrial loops and
the very large number of tuning methods proposed for it, a sign that tuning,
more than structure, is where the difficulty lies \citep{Astrom2001,Ang2005}.

\begin{figure*}[t]
\centering
\resizebox{\textwidth}{!}{\begin{tikzpicture}[
  font=\footnotesize,
  box/.style={draw, rounded corners=2pt, align=center, minimum height=9mm, text width=21mm, inner sep=2pt},
  pid/.style={box, fill=blue!6},
  lq/.style={box, fill=orange!8},
  our/.style={box, fill=green!10, very thick},
  arr/.style={-{Stealth[length=2mm]}, thick},
  lab/.style={font=\scriptsize\itshape, align=center}]
% PID branch
\node[pid] (p)   {P\\ \scriptsize Minorsky 1922};
\node[pid, right=5mm of p] (pi)  {PI\\ \scriptsize zero steady error};
\node[pid, right=5mm of pi] (pidn) {PID\\ \scriptsize Ziegler--Nichols 1942};
\node[pid, right=5mm of pidn] (pidf) {PIDF\\ \scriptsize filtered $D$, pole $N$};
\node[pid, right=5mm of pidf] (pids) {PID + setpoint\\ \scriptsize weights $b,c$};
\node[our, below=9mm of pids] (dof) {2-DOF PIDF\\ \scriptsize reads $v_o$ only};
\draw[arr] (p) -- (pi);
\draw[arr] (pi) -- (pidn);
\draw[arr] (pidn) -- (pidf);
\draw[arr] (pidf) -- (pids);
\draw[arr] (pids) -- (dof);
% LQ branch
\node[lq, below=25mm of p] (lqr) {LQR\\ \scriptsize Kalman 1960};
\node[lq, right=5mm of lqr] (lqg) {LQG\\ \scriptsize Athans 1971};
\node[lq, right=5mm of lqg] (lqi) {LQI\\ \scriptsize integral augmentation};
\node[our, right=5mm of lqi] (cmo) {CMO--LQI\\ \scriptsize constrained weights};
\draw[arr] (lqr) -- (lqg);
\draw[arr] (lqr.south) to[out=-20,in=-160,looseness=0.7] (lqi.south);
\draw[arr] (lqi) -- (cmo);
\draw[arr, very thick, green!40!black] (cmo.east) -| node[lab, pos=0.75, right, text=green!30!black] {exact map\\ Prop.~\ref*{prop:map}--\ref*{prop:struct}} (dof.south);
% legends
\node[lab, above=1mm of pidn] {output feedback, tuned by rules or search};
\node[lab, below=9mm of lqg] {state feedback, tuned by weights $Q,R$};
\end{tikzpicture}
}
\caption{Evolution of the two controller families used in this paper. Upper
branch: output-feedback PID-type laws, from the proportional controller to the
2-DOF PIDF. Lower branch: linear--quadratic state feedback, from LQR to LQI and
to the constrained weight tuning (CMO--LQI) proposed here. The two branches meet
through the exact realization of Propositions~\ref{prop:map} and
\ref{prop:struct}.}
\label{fig:evolution}
\end{figure*}

\subsubsection{From LQR to LQI}
The lower branch of Fig.~\ref{fig:evolution} starts with \citet{Kalman1960},
who posed the regulation of a linear system as the minimization of a quadratic
cost and showed that its solution is a constant state feedback computed from a
Riccati equation. The LQR comes with a certificate that no PID tuning rule
offers: for a single input, the return difference satisfies $|1+L(j\omega)|\ge1$,
which gives an infinite upward gain margin, a gain reduction margin of one half
and at least \SI{60}{\degree} of phase margin \citep{Kalman1964,Safonov1977,Anderson1990}.

The certificate assumes that every state is measured. When the states are
estimated by a Kalman filter, the controller becomes the LQG law
\citep{Athans1971}, and \citet{Doyle1978} showed with a counter-example that
LQG loops can have arbitrarily small margins; loop-transfer recovery restores
them only at the cost of observer bandwidth \citep{Doyle1979}. The second
limitation of the LQR is that a proportional state feedback does not remove the
steady error caused by a constant load or by a model error. The servomechanism
theory of \citet{Young1972} and \citet{Davison1976}, grounded on the internal
model principle of \citet{Francis1976}, solves this by adding the integral of
the tracking error to the state vector and solving the LQR for the augmented
plant. This linear--quadratic--integral (LQI) law is what the rest of the paper
uses.

Several variants of the LQI have been applied to converters. The weights have
been chosen by Bryson's rule \citep{Bryson1975}, by hand, or by an optimizer;
robust versions replace the Riccati equation by linear matrix inequalities that
hold over a polytope of operating points \citep{Olalla2009}, and closed-loop
poles can be confined to a prescribed region by the same means
\citep{Chilali1996};
early work already cast converter regulation as a general LQR problem
\citep{Leung1991}. Saturation of the duty ratio requires an anti-windup
mechanism, for which \citet{Kothare1994} gave a unified framework and
back-calculation remains the usual industrial choice \citep{Astrom1989}.
Other families exist and are strong on converters, for instance sliding-mode
control \citep{Tan2008}, hybrid model predictive control
\citep{Geyer2008,Liu2018}, and observer-based or intelligent nonlinear laws
\citep{Abdelmalek2020,Sitbon2020,Bouchama2020}, but they give up either the fixed switching
frequency or the low-order linear structure that a PID block can host. The LQI
is chosen here for three reasons that are specific to the Buck: the averaged
plant is second order, so the augmented state has dimension three and the
Riccati solution is cheap enough to be computed inside an optimizer; the
integral state gives zero steady error across a range in which the equilibrium
duty moves from 0.03 to 0.97; and, as shown in Section~\ref{sec:theory}, the
capacitor ESR places a zero in the output equation that makes the LQI exactly
realizable from the output voltage.

\subsubsection{Why a non-ideal Buck converter}
The Buck converter is often used as a benchmark because its ideal averaged
model is linear in the states and bilinear in the duty ratio
\citep{Middlebrook1976,Erickson2020}. For the question asked here the ideal
model is not enough. The ESR $r_C$ creates the output-network zero on which the
exact equivalence rests, so it cannot be dropped. The other parasitic elements
matter for the operating point. Each volt lost in $r_L$, $R_{on}$, $r_d$ or across
the diode threshold $V_f$ has to be made up by a larger duty ratio, and at a
\SI{1}{V} output these drops are no longer small compared with the output
itself: our model gives an efficiency of about \SI{61}{\percent} there, so
nearly \SI{40}{\percent} of the input power ends up in these elements
\citep{Kazimierczuk2015,Siddhartha2018}. The freewheeling diode adds a second
effect. It blocks reverse current, so at light load or during a deep
step-down the inductor current falls to zero inside the switching period and
the converter enters discontinuous conduction, where the small-signal model
loses one order \citep{Sun2001}. Our \SI{1}{V} operating point sits only
\SI{3}{\percent} above the conduction boundary at the nominal load and spends
\DCMone{} of its steady-state time in discontinuous conduction. Averaging itself
is an approximation whose validity depends on the ripple and the switching
frequency \citep{Krein1990,Maksimovic2001}, and in digital loops the sampled
ripple and the duty resolution can sustain limit cycles that no averaged model
predicts \citep{Peterchev2003,Peng2007,Buso2015}. A wide-range, non-ideal Buck
therefore combines saturation, a mode change and switching nonlinearities in a
plant that is still simple enough to analyse in closed form, which makes it a
fair test of any claim of equivalence.

\subsection{Research gap and motivation}\label{sec:gap}
The review leaves three gaps open. Each one motivates a design choice of this
paper, and we state them in pairs.

\subsubsection{Gap 1: the LQI--PID link is known only in nominal form}
The link between the LQI of a Buck converter and a filtered PID is known as a
transfer-function identity at one operating point. Whether it survives duty
saturation, anti-windup, discontinuous conduction and parameter mismatch has
not been shown, yet these are exactly the conditions a converter meets across a
wide range. Without that result, a PID tuned from an LQI is only an
approximation of it, and the LQ margins do not carry over.

\subsubsection{Motivation 1: why a PID equivalent, and why the 2-DOF PIDF}
If the LQI is better founded, one may ask why it should be turned into a PID at
all. There are three reasons. The LQI needs the inductor current, hence a
current sensor with its own bandwidth, cost and noise, whereas the output
voltage is measured on every converter. Commissioning engineers can inspect,
bound and retune three gains, and the PID block comes with validated
anti-windup and bumpless transfer \citep{Astrom2006}. And tuning a PID directly
in gain space, by rules or by an optimizer, gives no stability or robustness
guarantee for the point it returns \citep{Ang2005}. An exact map from the LQI to
a PID-type law removes that last objection: every candidate of a search becomes
a controller that has the LQR margins and still runs on a standard block.

The map needs the extra parameters of the 2-DOF PIDF. The LQI feeds back the
capacitor voltage and the inductor current; replacing them by the measured
output requires a filtered derivative whose pole sits on the ESR zero, and that
pole is the $N$ of the PIDF. The LQI also sees the reference only through its
integral state, so the equivalent PIDF must have $b=1$ and $c=0$. Without
setpoint weights a PIDF can match the disturbance response of the LQI but not
its command response (Proposition~\ref{prop:map}). A PI or a PID lacks the
parameters for either condition, which is why they appear here only as
comparators. Table~\ref{tab:ctrl} lists all the controllers of the study.

\begin{table*}[t]
\centering
\caption{Controllers used in this study.}
\label{tab:ctrl}
\footnotesize
\resizebox{\textwidth}{!}{%
\begin{tabular}{@{}llllll@{}}
\toprule
Controller & Measured signals & Tuning parameters & Certificate & Role here & Key references\\
\midrule
PI & $v_o$ & $K_p,K_i$ & none & comparator & \citet{Ziegler1942}\\
PID & $v_o$ & $K_p,K_i,K_d$ & none & comparator & \citet{Astrom2006}\\
PIDF & $v_o$ & $K_p,K_i,K_d,N$ & none & comparator & \citet{Ang2005}\\
2-DOF PIDF & $v_o$ & $K_p,K_i,K_d,N,b,c$ & none if tuned directly & realization of the proposal & \citet{Horowitz1963}\\
LQR & $i_L,v_o$ & $Q,R$ & LQ margins & comparator & \citet{Kalman1960}\\
LQG & $v_o$ (+ observer) & $Q,R$, noise covariances & none & comparator & \citet{Athans1971}\\
LQI & $i_L,v_o$ & $Q,R$ & LQ margins, zero steady error & comparator & \citet{Young1972}\\
CMO--LQI--PIDF & $v_o$ & 4 log-excursions & LQ margins, D-region & proposal & this work\\
\bottomrule
\end{tabular}}
\end{table*}

\subsubsection{Gap 2: metaheuristics tune gains, not specifications}
Metaheuristic tuning of converter controllers searches gain spaces
\citep{Izci2022,Nanyan2024,Siano2014,ValenciaRivera2024}. A good result
therefore carries no guarantee, the size of the search box is a guess, and the
objective is usually a tracking index rather than the list of limits a
specification actually contains. Searching the LQI weights instead has not been
compared with a direct gain search under the same budget.

\subsubsection{Motivation 2: CMO--LQI, constrained multi-criteria optimization of the LQI weights}
Bryson's rule gives a sensible starting point, but its weights are not optimal
with respect to overshoot, device current or loss. We call CMO--LQI
(constrained multi-criteria optimization of the LQI) the problem of choosing the
weights so that the resulting LQI, realized as a 2-DOF PIDF, minimizes an
aggregate of seven performance criteria over six scenarios under hard
engineering constraints. \emph{Multi-criteria} is meant literally, not in the
Pareto sense: the criteria are normalized and averaged into one scalar, and the
constraints are handled by feasibility rules \citep{Deb2000}. A weighted scalar
index does not describe a trade-off front \citep{RodriguezMolina2020}, and we do
not claim one. The search variables are the logarithms of the Bryson excursions
and of the integral bandwidth, four numbers with a physical scale, rather than
the PIDF gains.

\subsubsection{Gap 3: candidates are accepted on the model that selected them}
Most tuning studies accept a candidate on the averaged model that drives the
search. In an earlier version of the present protocol, four of six winners
accepted in this way failed on the switched model, because of switching limit
cycles that the averaged model cannot represent. The same studies often
introduce a new operator without a matched comparison against existing ones,
so it is hard to tell whether a result comes from the algorithm or from the
problem.

\subsubsection{Motivation 3: why metaheuristics, and why refineBO}
The CMO--LQI problem has only four variables, but its objective and its
constraints come from simulations with saturation, mode changes and switching,
so no gradient is available and the feasible set may be disconnected.
Population-based metaheuristics are the usual tool for such black-box problems.
Two results from the optimization literature shape the way we use them. The
no-free-lunch theorems state that no algorithm dominates on all problems
\citep{Wolpert1997}, and metaphor-based variants are often published without
evidence that they differ from existing methods \citep{Sorensen2015}. We
therefore do not propose a new metaheuristic. We take three classical operators
with well-documented behaviour, particle swarm optimization
\citep{Kennedy1995,Poli2007}, the genetic algorithm \citep{Holland1975} and the
artificial bee colony \citep{Karaboga2007,Karaboga2008}, and three recent
variants built on the grey wolf optimizer and on an adaptive bee colony
\citep{Mirjalili2014,NadimiShahraki2021,Hou2022,Alsamia2024}, and ask whether the
answer depends on the operator. Constraints are handled by Deb's feasibility
ordering, which avoids penalty factors \citep{Deb2000,MezuraMontes2011}; the
runs are compared on matched seeds with non-parametric tests
\citep{Derrac2011,Demsar2006}; and three of the 33 constraints are measured on
a short switched simulation inside every evaluation, which answers the
acceptance problem of Gap 3.

A global operator spends most of its budget exploring. Once a promising region
is found, a local model-based stage is more economical, which is the idea of
memetic and hybrid schemes \citep{Neri2012}. For expensive black-box functions,
Bayesian optimization fits a Gaussian-process surrogate to the evaluated points
and chooses the next one by maximizing an acquisition function such as the
expected improvement \citep{Jones1998,Rasmussen2006,Snoek2012,Shahriari2016}. Unknown
constraints are handled by multiplying the acquisition by the modelled
probability of feasibility \citep{Gardner2014}, and a trust region keeps the
surrogate local and reliable \citep{Eriksson2019}. Bayesian optimization has
already been used to tune industrial controllers \citep{NeumannBrosig2020} and
digital PID loops \citep{Coutinho2023}. The refineBO stage used here combines
these three elements: it starts from the history of each global run, models the
objective and all 33 constraints with Gaussian processes, and searches a
shrinking region around the best feasible records.

Table~\ref{tab:gap} positions the work against representative Buck control and
tuning studies. The motivation, in short, is to close the three gaps on a plant
that is fully specified, so that every number in the paper can be recomputed.

\begin{table*}[t]
\centering
\caption{Positioning with respect to representative Buck control and tuning studies. ``Not reported'' refers to the stated scope of each work only.}
\label{tab:gap}
\footnotesize
\resizebox{\textwidth}{!}{%
\begin{tabular}{@{}p{3.0cm}p{3.3cm}p{3.6cm}p{6.4cm}@{}}
\toprule
Study & Plant evidence & Controller / search & Relative to the present question\\
\midrule
\citet{Siddhartha2018} & non-ideal Buck, dSPACE & IMC-PID, direct & no repeated solver comparison, no switched energy audit\\
\citet{Izci2022} & averaged Buck & PID, AEO + Nelder--Mead & gain-space search, no structural certificate, no switched admissibility\\
\citet{Nanyan2024} & averaged Buck & PID, improved SCA & gain-space search, single fidelity\\
\citet{Siano2014} & Buck prototype & fuzzy, MO-PSO & no equal-budget solver comparison\\
\citet{Abdelmalek2020} & Buck, experimental & observer-based nonlinear, PSO & tracking/robustness focus, not a matched co-design\\
\citet{Olalla2009} & PWM converters, polytopic model & robust LQR & polytopic certificate, no PID realization, no search over weights\\
\citet{Coutinho2023} & multiloop process & PID, Bayesian optimization & not a power converter\\
Present work & audited non-ideal Buck, 1--46 V, switched replica with energy balance & LQI weights $\to$ exact 2-DOF PIDF; 6 global operators + refineBO; 6 families & structural equivalence theorem; 33 inequalities (3 measured on the switched model); switched validation vs.\ PI, PID, PIDF, LQR, LQG, LQI\\
\bottomrule
\end{tabular}}
\end{table*}

\subsection{Assumptions}\label{sec:assump}
The analysis rests on the following assumptions, each of which is tested or
bounded later in the paper.
\begin{itemize}
\item The input voltage is constant, $V_{in}=\SI{48}{V}$; line regulation is not
  part of the specification.
\item The output network $(R,r_C,C)$ is known at design time. Load steps are
  treated as unmeasured disturbances, and the error they cause in the
  equivalence is given in closed form.
\item The PIDF measures the output voltage only. The inductor current is
  measured only by the full-state comparators and by the state-feedback twin of
  the proposal.
\item Control is sampled synchronously with the \SI{50}{kHz} carrier, with a
  one-sample computational delay; ADC quantization and the finite resolution of
  the PWM are not modelled.
\item Losses are the conduction and snubber losses of the declared model
  boundary; switching, reverse-recovery, magnetic and thermal losses are
  outside it.
\item Component tolerances were not supplied. Robustness claims are made on
  the nominal operating polytope only.
\end{itemize}

\subsection{Challenges}\label{sec:challenges}
Three difficulties had to be solved. The first is theoretical: showing that
the equivalence holds along trajectories, not only in the frequency domain,
and identifying exactly which disturbances break it. The second is numerical.
A zero-order-hold discretization of the derivative filter, the usual choice,
introduces a lag that separates the sampled PIDF from the sampled LQI by about
\EqSampZOH{} during a wide slew, so the discretization had to be derived from
the observer interpretation of the filter. The third concerns the search. The
feasible set is defined by 33 inequalities, three of which depend on switching
phenomena; checking them on the switched model is expensive, while ignoring
them lets through candidates that are unacceptable on the real power stage.
The screen had to be short enough to run inside every evaluation and long
enough to excite the coexisting limit cycles.

\subsection{Contributions}\label{sec:contrib}
The contributions of the paper are the following.
\begin{enumerate}
\item A formula-first chain from the non-ideal converter equations to the
  exact output-feedback realization of any LQI law as a 2-DOF PIDF, with every
  analytical quantity (equilibria, transfer functions, Riccati gain, LQR
  margins, a common D-region certificate over the operating polytope)
  instantiated and independently recomputed.
\item A structural equivalence theorem (Proposition~\ref{prop:struct}): the
  PIDF filter state equals the capacitor voltage along every trajectory. Its
  domain of validity is stated, the gap to the full-state law is given in
  closed form, and a first-order-hold discretization reduces the sampled gap
  from \EqSampZOH{} to \EqSampFOH{}.
\item A constrained multi-criteria formulation (CMO--LQI) in the space of LQI
  weights, with 33 inequalities in engineering units, three of which are
  measured on the switched model during the search, together with the evidence
  from two pilots of why an analytical ripple bound or a steady switched screen
  is not sufficient.
\item A matched, equal-budget comparison of seven algorithms (six global
  operators followed by refineBO) across six controller families. It shows that
  the LQI coordinates do not speed up the search but confine it, at no cost in
  optimality, to the LQ-optimal set on which the gain-space optimum also lies.
\item A switched-model validation on the \SI{24}{}$\to$\SI{1}{}$\to$\SI{46}{V}
  mission against PI, PID, PIDF, LQR, LQG and LQI, reported specification by
  specification.
\end{enumerate}

\subsection{Paper organization}\label{sec:org}
Section~\ref{sec:model} derives the non-ideal model and describes the switched
replica of the reference Simulink model. Section~\ref{sec:theory} states the
LQI synthesis and proves the nominal and structural equivalences.
Section~\ref{sec:problem} defines the CMO--LQI problem and the seven algorithms.
Section~\ref{sec:results} reports the verification of the equivalence, the
algorithm and family comparisons, and the switched validation, and discusses
the limitations. Section~\ref{sec:concl} concludes. The code, including a MATLAB
S-function that runs the same control layer inside the reference Simulink
model, is released with the paper.
